% Level 1, Homework 9 -- Mini Math Quiz   (input by homework.tex)
\homeworkweek{9}{Mini Math Quiz}{15--20 minutes. Do Core first; Challenge and Creative are optional.}
\namedateline

\begin{core}[Mini Math Quiz (about 10 minutes)]
Open \texttt{week09/hw-math-quiz/starter.py}. No turtle this time: everything happens in the Shell.
\begin{enumerate}[leftmargin=1.4em,itemsep=1pt]
  \item Press Run and answer question 1 in the Shell. Read how it works: \texttt{input}, then \texttt{int}, then \texttt{if}/\texttt{else}, and \py|score = score + 1| when the answer is right.
  \item Add question 2 the same way: \enquote{What is 9 + 15?} (the answer is 24).
  \item At the end, print the score: \py|print("You scored", score, "out of 2")|.
  \item Test it twice: once with right answers, once with wrong ones. Then run \texttt{check\_my\_work.py}.
\end{enumerate}
\textbf{No computer at home?} Write question 2 here (copy the pattern of question 1):
\answerlines{6}
\end{core}

\begin{challenge}[Perfect score]
Add a third question of your own. Change \enquote{out of 2} to \enquote{out of 3}. Then, at the very end, add an \texttt{if} that prints a special message when \texttt{score == 3}. Remember: two equals signs ask a question!
\end{challenge}

\begin{creative}[Choose your path]
Write a tiny adventure: the program asks \enquote{You find a door and a tunnel. Which do you choose?} and uses \texttt{if}/\texttt{elif}/\texttt{else} to print a different ending for \texttt{door}, \texttt{tunnel} and anything else. Plan it here first:
\answerlines{4}
Remember the rule: our programs ask about \textbf{choices}, never about who you are (no names, addresses, schools or passwords).
\end{creative}

\begin{selfcheck}
  \item My quiz asks 2 questions and says \texttt{Correct!} for right answers.
  \item The last line shows my score.
  \item I know why we need \texttt{int(...)} around \texttt{input(...)} for numbers.
  \item I know the difference between \texttt{=} (put in the box) and \texttt{==} (are they the same?).
\end{selfcheck}
